\begin{table}[H]
\centering
\caption{Standardized Effect Sizes for Main Outcomes}
\label{tab:sde}
\begin{adjustbox}{max width=\textwidth}
\begin{tabular}{llccccc}
\toprule
Outcome & Specification & $\hat{\beta}$ & SD($Y$) & SDE & SE(SDE) & Classification \\
\midrule
\multicolumn{7}{l}{\textit{Panel A: Pooled}} \\
OKP costs/insured & Std.\ dose $\times$ Post-2008 & 67.81 & 796 & 0.0852 & 0.0221 & Moderate positive \\
\midrule
\multicolumn{7}{l}{\textit{Panel B: Heterogeneous}} \\
~~German-speaking & DI rate $\times$ Post-2008 & 6.947 & 756 & 0.0947 & 0.0251 & Moderate positive \\
~~French/Italian & DI rate $\times$ Post-2008 & 5.276 & 754 & 0.0721 & 0.0366 & Moderate positive \\
\bottomrule
\end{tabular}
\end{adjustbox}
\par\vspace{0.3em}
{\footnotesize \textit{Notes:} \textbf{Country:} Switzerland. \textbf{Research question:} Whether Switzerland's disability insurance reforms (5th revision 2008, 6a revision 2012), which shifted from compensating disability to preventing it through early intervention and reintegration, affected mandatory health insurance (OKP) per-capita costs. \textbf{Policy mechanism:} The IV reforms introduced early intervention measures, integration programs, and vocational rehabilitation to keep individuals with health limitations in the labor force rather than granting disability pensions, potentially altering their use of health services through changed health trajectories or substitution from disability benefits to medical treatment. \textbf{Outcome definition:} OKP gross benefits per insured person (Bruttoleistungen pro Versicherten), measuring total mandatory health insurance spending per capita in Swiss francs. \textbf{Treatment:} Continuous; DI pension rate per 1,000 cantonal population in 2009 (standardized, mean zero, SD one), measuring pre-reform disability burden and thus exposure to the reforms. \textbf{Data:} BAG Dashboard Krankenversicherung OKP (health costs, 1997--2024) merged with BFS PXWeb IV statistics (disability, 2009--2024); 26 cantons, 2000--2022, 598 canton-year observations. \textbf{Method:} Dose-response difference-in-differences with canton and year fixed effects; standard errors clustered at canton level. \textbf{Sample:} 26 Swiss cantons, excluding 2023--2024 due to potentially incomplete reporting; the analysis exploits cross-canton variation in disability insurance burden as a continuous treatment dose. SDE $= \hat{\beta} / \text{SD}(Y)$ where SD($Y$) is the unconditional standard deviation of OKP costs per insured. Classification refers to magnitude, not statistical significance: Large ($|$SDE$| > 0.15$), Moderate ($0.05$--$0.15$), Small ($0.005$--$0.05$), Null ($< 0.005$).}
\end{table}
